% ===========================================================================
%  CAPÍTULO 12 — CONCLUSIONES
% ===========================================================================

\section*{12. Conclusiones}
\label{sec:conclusiones}

\subsection{12.1. Cumplimiento de objetivos}
\label{subsec:cumplimiento_objetivos}

El desarrollo de la plataforma \textit{LigaMaster} ha concluido con
éxito, logrando el objetivo general propuesto al inicio del trabajo:
diseñar e implementar una plataforma web de apoyo a la toma de
decisiones en el Fantasy Football mediante análisis predictivo y
explicabilidad.

Se ha logrado una integración completa entre un motor de
\textit{Machine Learning} desarrollado en Python, un \textit{backend}
robusto con Django REST Framework y un \textit{frontend} interactivo y
desacoplado creado con React. Los objetivos funcionales se han cubierto
en su totalidad, permitiendo la visualización de estadísticas, la gestión
de plantillas y la predicción de rendimiento basada en explicabilidad
(XAI). Del mismo modo, los requisitos no funcionales de rendimiento
(RNF1), calidad del dato (RNF3) y seguridad (RNF5) han sido
verificados mediante la suite de pruebas automatizadas, las métricas de
cobertura y los análisis estáticos con SonarQube descritos en el
capítulo anterior.

% ---------------------------------------------------------------------------
\subsection{12.2. Limitaciones encontradas}
\label{subsec:limitaciones}

A lo largo del proyecto se han identificado ciertas limitaciones
inherentes a la obtención y al procesamiento de los datos, así como a
las decisiones de diseño adoptadas bajo las restricciones propias de un
desarrollo individual. Las más relevantes se describen a continuación:

\begin{itemize}
  \item \textbf{Ausencia de APIs oficiales para el mercado.} La falta
        de \textit{endpoints} públicos en la API oficial de LaLiga
        Fantasy impidió extraer e integrar los valores económicos
        dinámicos de los jugadores en tiempo real de forma automatizada.
        Esto ha limitado la capacidad de desarrollar un módulo
        financiero completo en la plataforma, funcionalidad que habría
        enriquecido considerablemente la propuesta y habría permitido
        correlacionar el precio de mercado con la predicción de
        rendimiento.

  \item \textbf{Bloqueos de IP durante el \textit{scraping}.} Debido
        a los mecanismos anti-\textit{scraping} activos en las fuentes
        consultadas, fue necesario descargar manualmente el HTML de
        todos los partidos históricos, proceso que resultó costoso en
        tiempo y difícilmente automatizable. Esta restricción implica
        que la actualización del \textit{dataset} requiere intervención
        manual periódica, lo que supone un cuello de botella operativo
        en un entorno de producción continua.

  \item \textbf{Capacidad de carga del sistema.} Las pruebas de
        rendimiento con Locust evidenciaron que, en su configuración
        actual, el sistema comienza a degradarse significativamente a
        partir de los 20--30 usuarios concurrentes, con tasas de error
        del 14,5\,\% a 100 usuarios. El diagnóstico apunta al agotamiento
        del número máximo de conexiones simultáneas de PostgreSQL como
        causa principal, lo que constituye una limitación relevante de
        cara a un despliegue en producción con tráfico real.
\end{itemize}

% ---------------------------------------------------------------------------
\subsection{12.3. Líneas de trabajo futuro}
\label{subsec:trabajo_futuro}

El sistema desarrollado posee una arquitectura escalable que abre la
puerta a diversas mejoras y nuevas vías de desarrollo. A continuación
se presentan las líneas de trabajo consideradas de mayor impacto, tanto
a nivel técnico como de producto:

\begin{itemize}

  \item \textbf{Optimización del rendimiento y escalabilidad.}
        El cuello de botella identificado en las pruebas de carga puede
        resolverse mediante la incorporación de un \textit{connection
        pooler} como PgBouncer entre Django y PostgreSQL, lo que
        permitiría gestionar un número elevado de peticiones concurrentes
        sin saturar el límite de conexiones de la base de datos.
        Complementariamente, ampliar el uso de Redis como capa de caché
        para cubrir los \textit{endpoints} de mayor coste computacional
        —especialmente los de predicción y clasificación— reduciría la
        presión directa sobre el gestor relacional. Estas dos mejoras,
        de bajo coste de implementación, permitirían escalar el sistema
        a un orden de magnitud superior sin modificar la arquitectura
        existente.

  \item \textbf{Consolidación y refactorización de la capa API.}
        Durante el desarrollo del proyecto coexistieron dos capas de
        acceso a datos: la API principal estructurada en el módulo
        \texttt{main/api/} y un conjunto de adaptadores de compatibilidad
        en \texttt{main/drf\_views.py}, introducidos para mantener rutas
        equivalentes durante una migración interna. Esta situación genera
        deuda técnica al duplicar puntos de entrada, dificultar el
        mantenimiento y aumentar la superficie de posibles
        inconsistencias. Como trabajo futuro prioritario se propone
        unificar ambas capas bajo una única interfaz coherente,
        eliminando los adaptadores de compatibilidad y normalizando el
        contrato de todos los \textit{endpoints} bajo un versionado
        explícito (\texttt{/api/v1/}).

  \item \textbf{Automatización del \textit{pipeline} de reentrenamiento.}
        En su estado actual, los modelos de Machine Learning son
        entrenados de forma manual fuera de línea y serializados como
        artefactos \texttt{.pkl} estáticos. Este enfoque, adecuado para
        el alcance del TFG, no es sostenible en producción: el
        rendimiento de los modelos se degrada con el tiempo a medida
        que el contexto de la competición evoluciona, y la incorporación
        de nuevas jornadas requiere intervención manual. Una línea de
        trabajo natural sería implementar un \textit{pipeline} de
        reentrenamiento periódico y automatizado —por ejemplo, al cierre
        de cada jornada— utilizando una cola de tareas asíncronas como
        Celery. Este enfoque se enmarca dentro de las prácticas de
        MLOps y garantizaría que los modelos desplegados reflejen
        siempre el estado más reciente de la competición.

  \item \textbf{Predicción de titularidades.} Integrar un modelo
        secundario que evalúe la probabilidad de que un jugador parta
        en el once inicial. Mediante el análisis de rotaciones tácticas,
        el histórico de minutos jugados, la acumulación de tarjetas y
        los partes médicos disponibles, se refinaría notablemente la
        precisión de las puntuaciones esperadas por jornada. Actualmente
        el sistema asume la disponibilidad plena del jugador, lo que
        constituye una simplificación relevante que limita la fiabilidad
        de las predicciones en jornadas con alta rotación.

  \item \textbf{Desarrollo de una aplicación móvil nativa.}
        Aprovechando que el \textit{backend} está completamente
        desacoplado y expuesto a través de una API REST con DRF, el
        siguiente paso natural sería implementar una aplicación nativa
        para dispositivos iOS y Android. Esto mejoraría la experiencia
        de usuario en el entorno donde la mayoría de personas acceden
        a los juegos de Fantasy Football y permitiría habilitar
        notificaciones \textit{push} en tiempo real, complementando el
        sistema de WebSockets ya implementado.

\end{itemize}

% ---------------------------------------------------------------------------
\subsection{12.4. Valoración personal}
\label{subsec:valoracion_personal}

La realización de este Trabajo de Fin de Grado ha supuesto un verdadero
reto que me ha permitido consolidar y poner en práctica los conocimientos
adquiridos a lo largo de la titulación. A través de este proyecto he
gestionado todas las etapas del ciclo de vida del software: desde la
captura de requisitos y el diseño de procesos ETL, hasta la
implementación de modelos de aprendizaje automático, el desarrollo
\textit{full-stack} y el despliegue final.

Si bien estas tareas han sido recurrentes durante mi formación, la
ejecución de este proyecto ha supuesto un hito personal al asumir, por
primera vez, la autonomía y la responsabilidad total sobre un sistema
complejo de principio a fin.

Especialmente valioso ha resultado el aprendizaje autodidacta en técnicas
de Explainable AI (XAI), lo que me ha dotado de una perspectiva diferente
y realista sobre las exigencias actuales en el desarrollo de software y
en el mundo de la inteligencia artificial. Asimismo, el proyecto ha
permitido profundizar en el tratamiento de datos a gran escala y en los
fundamentos prácticos de la ciencia de datos aplicada a un dominio real.
En conjunto, considero que \textit{LigaMaster} refleja fielmente el
perfil técnico que una titulación en Ingeniería del Software debe
desarrollar: no solo la capacidad de construir sistemas funcionales, sino
de razonar sobre sus limitaciones, documentar sus decisiones y diseñar
una hoja de ruta que permita su evolución responsable.